
\documentclass{article}
\usepackage{colortbl}
\usepackage{makecell}
\usepackage{multirow}
\usepackage{supertabular}

\begin{document}

\newcounter{utterance}

\centering \large Interaction Transcript for game `eqbench', experiment `eq\_bench', episode 11 with reflexion{-}Qwen3.5{-}4B.
\vspace{24pt}

{ \footnotesize  \setcounter{utterance}{1}
\setlength{\tabcolsep}{0pt}
\begin{supertabular}{c@{$\;$}|p{.15\linewidth}@{}p{.15\linewidth}p{.15\linewidth}p{.15\linewidth}p{.15\linewidth}p{.15\linewidth}}
    \# & \multicolumn{2}{c}{Player} && \multicolumn{2}{c}{Game Master} \\
    \hline

    \theutterance \stepcounter{utterance}  
    & & & \multicolumn{4}{p{0.6\linewidth}}{
        \cellcolor[rgb]{0.9,0.9,0.9}{
            \makecell[{{p{\linewidth}}}]{
                \texttt{\tiny{[P1$\langle$GM]}}
                \texttt{You have attempted this same task before and did not succeed. These are the plans you wrote after those attempts.} \\
\texttt{Trial \#1: New plan: I failed because I assigned a score of 1 to "Admiration," which is likely too high given Harold's dismissive tone ("It's not like the world is going to end") and his questioning of Lorelei's candidate ("And you think your candidate has all the answers?"). Harold is skeptical and frustrated, not impressed. I should have set Admiration to 0. Additionally, I need to ensure that if I assign a score > 0 to an emotion, I am confident it is the primary or secondary feeling based on the text. Given Harold's irritation at being called "worked up" and his skepticism, Irritation should be the highest score, and Disbelief might be moderate. I will re{-}evaluate the scores to reflect Harold's frustration and lack of respect for Lorelei's perspective, ensuring Admiration is 0.} \\
\\ 
\texttt{Disbelief: 4} \\
\texttt{Irritation: 9} \\
\texttt{Admiration: 0} \\
\texttt{Indifference: 0} \\
\\ 
\\ 
\\ 
\texttt{Your task is to predict the likely emotional responses of a character in this dialogue:} \\
\\ 
\texttt{Harold: I don't understand why you're so worked up about this election. It's not like the world is going to end if your candidate doesn't win.} \\
\texttt{Lorelei: That's not the point, Harold. It's about what this election represents. It's about the future of our planet, our society.} \\
\texttt{Harold: And you think your candidate has all the answers? } \\
\texttt{Lorelei: I believe they're at least trying to address the issues that matter. Climate change, social justice, healthcare...} \\
\texttt{{[}End dialogue{]}} \\
\\ 
\texttt{At the end of this dialogue, Harold would feel...} \\
\texttt{Disbelief} \\
\texttt{Irritation} \\
\texttt{Admiration} \\
\texttt{Indifference} \\
\\ 
\texttt{Give each of these possible emotions a score from 0{-}10 for the relative intensity that they are likely to be feeling each.} \\
\\ 
\texttt{You must output in the following format, including headings (of course, you should give your own scores), with no additional commentary:} \\
\\ 
\texttt{Disbelief: <score>} \\
\texttt{Irritation: <score>} \\
\texttt{Admiration: <score>} \\
\texttt{Indifference: <score>} \\
\\ 
\\ 
\texttt{{[}End of answer{]}} \\
\\ 
\texttt{Remember: zero is a valid score, meaning they are likely not feeling that emotion. You must score at least one emotion > 0.} \\
\\ 
\texttt{Your answer:} \\
            }
        }
    }
     \\ \\

    \theutterance \stepcounter{utterance}  
    & \multicolumn{4}{p{0.6\linewidth}}{
        \cellcolor[rgb]{0.561,0.737,0.561}{
            \makecell[{{p{\linewidth}}}]{
                \texttt{\tiny{[P1$\rangle$GM]}}
                \texttt{Disbelief: 3} \\
\texttt{Irritation: 9} \\
\texttt{Admiration: 0} \\
\texttt{Indifference: 0} \\
            }
        }
    }
    & & \\ \\

    \theutterance \stepcounter{utterance}  
    & & & \multicolumn{2}{p{0.3\linewidth}}{
        \cellcolor[rgb]{0.9,0.9,0.9}{
            \makecell[{{p{\linewidth}}}]{
                \texttt{\tiny{[GM$|$GM]}}
                \texttt{Disbelief: 3} \\
\texttt{Irritation: 9} \\
\texttt{Admiration: 0} \\
\texttt{Indifference: 0} \\
            }
        }
    }
    & & \\ \\

    \theutterance \stepcounter{utterance}  
    & & & \multicolumn{2}{p{0.3\linewidth}}{
        \cellcolor[rgb]{0.9,0.9,0.9}{
            \makecell[{{p{\linewidth}}}]{
                \texttt{\tiny{[GM$|$GM]}}
                \texttt{Disbelief: 4} \\
\texttt{Irritation: 6} \\
\texttt{Admiration: 0} \\
\texttt{Indifference: 0} \\
            }
        }
    }
    & & \\ \\

    \theutterance \stepcounter{utterance}  
    & & & \multicolumn{2}{p{0.3\linewidth}}{
        \cellcolor[rgb]{0.9,0.9,0.9}{
            \makecell[{{p{\linewidth}}}]{
                \texttt{\tiny{[GM$|$GM]}}
                \texttt{game\_result = LOSE} \\
            }
        }
    }
    & & \\ \\

    \theutterance \stepcounter{utterance}  
    & & \multicolumn{4}{p{0.6\linewidth}}{
        \cellcolor[rgb]{0.95,0.95,0.95}{
            \makecell[{{p{\linewidth}}}]{
                \texttt{\tiny{[P1$\leftrightarrow$P]}}
                \texttt{Post{-}episode Reflexion (not sent to GM):} \\
\texttt{New plan: I failed because I assigned a score of 4 to "Disbelief," which is too high given the context. Harold's statement "I don't understand why you're so worked up" and his follow{-}up "It's not like the world is going to end" indicate he is dismissing Lorelei's passion as irrational or excessive, not that he genuinely doubts the election's outcome or the candidate's existence. His tone is more about annoyance at her intensity than a lack of belief in the event itself. Therefore, Disbelief should be 0. Irritation remains the dominant emotion due to his dismissive language and questioning of her candidate's competence. I will adjust Disbelief to 0 and keep Irritation high to accurately reflect his frustration with her emotional state rather than a lack of belief in the election.} \\
\\ 
\texttt{Disbelief: 0} \\
\texttt{Irritation: 9} \\
\texttt{Admiration: 0} \\
\texttt{Indifference: 0} \\
            }
        }
    }
    & \\ \\

\end{supertabular}
}

\end{document}
